\AllagmaClaim{directional-effects}
Figure~\ref{fig:effects} retains the original paired-effect figure. The
complete cell estimates and descriptive uncertainty are generated directly
from the retained analysis in Table~\ref{tab:cells}. The two intervals
entirely below zero are the moons, constant-rate, 5k comparisons. The
positive cell means at some completed cosine endpoints are small and have
intervals that include zero; they do not establish harm or equality.

\AllagmaFigure{effects}
\AllagmaTable{cells}

\AllagmaClaim{schedule-pattern}
Table~\ref{tab:schedule} gives the prespecified schedule contrasts. The
two intervals entirely above zero are again moons at 5k. The remaining
schedule contrasts have wider uncertainty and include zero. The repeated
positive direction is descriptive evidence within this design, not eight
independent discoveries or a mechanism test.

\AllagmaClaim{absolute-quality}
This distinction prevents a reduced relative EMA benefit from being
mistaken for poorer absolute EMA performance. The absolute state-level
estimates are retained in the ancillary analysis.

\AllagmaClaim{duration-uncertainty}
The duration and interaction estimates appear in the appendix. All their
intervals contain zero, which leaves substantial uncertainty rather than
demonstrating equivalent effects across durations.

\AllagmaClaim{coverage-ceiling}
The ceiling limits the diagnostic's ability to distinguish these model
states. SW1, component counts and inlier fractions remain separate retained
measurements.

\AllagmaTable{schedule}
